% 最优潮流总览与公共接口层
% 本章文件由 \input 引入，不含导言区。
\section{最优潮流总览与公共接口层}\label{sec:overview}

本章给出平台最优潮流（OPF）模块的全景与公共接口层：求解器家族的组成与
数据流、Native AC 与交直流混合引擎的对比、公共选项 \fld{ACOPFOptions}/\fld{DCOPFOptions}
的全部字段、结果类型 \fld{ACOPFResult}/\fld{DCOPFResult} 的全部字段及其诚实
口径（\fld{model\_scope}、\fld{ValidityFlags}、\fld{model\_limitations}）、
问题描述 \fld{OPFProblem} 与求解器抽象 \fld{IOPFSolver}、嵌入式 Ipopt 后端的
编译门控与运行时探测、OPF 在平台中的主要调用方，以及求解器能力的运行时查询。
各引擎的内部数学（建模层、IPM 求解核、RPO 离散搜索、三相混合 OPF）
不在本章展开，分别见第\ref{sec:acopfnative}、\ref{sec:parityform}、
\ref{sec:parityipm}、\ref{sec:dcopf}、\ref{sec:rpo}、\ref{sec:threephase}、
节；验证与校核通道见第\ref{sec:verification}节。AML 建模层属于独立的
\srcpath{power\_models} 模块，其公式见
\srcpath{docs/modules/power\_models/power\_models\_manual.tex}。

除特别注明外，本章源码范围是
\srcpath{include/hacdcpf/optimal\_power\_flow/} 目录下的公共头文件、
对外门面 \srcpath{include/hacdcpf/api/hacdcpf.hpp} 与
\srcpath{src/api/hacdcpf.cpp}、以及三个求解器实现文件中的路径分派代码。
单位制遵循全平台约定：MW$\cdot$kV$\cdot$MVA，$S_{\mathrm{base}}=100$~MVA，
内部方程一律标幺值；成本类量的币种在代码中未固定，参数表单位列填“—”。

\subsection{求解器家族全景与数据流}\label{subsec:overview-family}

\subsubsection*{家族成员与入口}

平台 OPF 家族由六个求解/建模组件构成，全部挂在命名空间
\srcpath{hacdcpf::opf}（三相混合 OPF 为 \srcpath{hacdcpf::opf::phase\_hybrid}）
之下：

\begin{itemize}
  \item \textbf{AC OPF（纯交流）}：非线性交流 OPF，统一入口
        \srcpath{opf::solve\_ac\_opf}
        （\srcpath{include/hacdcpf/optimal\_power\_flow/ac\_opf\_solver.hpp:solve\_ac\_opf}；
        实现 \srcpath{src/optimal\_power\_flow/ac\_opf.cpp:solve\_ac\_opf}）。
        纯交流算例默认走 Native 紧致路径（第\ref{sec:acopfnative}节）；
        后端选择规则见第\ref{subsec:overview-engines}节。
  \item \textbf{交直流混合 OPF}：含 DC/VSC/LCC 组件的算例由同一入口
        \srcpath{opf::solve\_ac\_opf} 在 Auto 后端下一律分派到 Parity 全空间
        建模（第\ref{sec:parityform}节）+ Parity IPM 或嵌入式 Ipopt 求解核
        （第\ref{sec:parityipm}节）；显式选择不兼容后端时拒绝求解而非静默
        丢弃直流网络（见第\ref{sec:verification}节"混合系统禁 AC-only 回退"）。
  \item \textbf{DC OPF}：无损线性化 LP/QP（纯交流口径，不建直流电网），入口
        \srcpath{opf::solve\_dc\_opf}
        （\srcpath{include/hacdcpf/optimal\_power\_flow/dc\_opf\_solver.hpp:solve\_dc\_opf}；
        实现 \srcpath{src/optimal\_power\_flow/dc\_opf.cpp:solve\_dc\_opf}），另提供
        独立于求解的可行性复核 \srcpath{check\_dc\_opf\_feasibility}
        （同头文件:22--25）。详见第\ref{sec:dcopf}节。
  \item \textbf{RPO}：无功优化 MINLP（连续内层 NLP + 离散分接头/并联补偿
        搜索），入口 \srcpath{opf::solve\_rpo}
        （\srcpath{src/optimal\_power\_flow/reactive\_power\_opt.cpp:solve\_rpo}；
        选项 \fld{RPOOptions} 定义于
        \srcpath{include/hacdcpf/optimal\_power\_flow/reactive\_power\_opt.hpp:RPOOptions}）。
        详见第\ref{sec:rpo}节。
  \item \textbf{三相混合 OPF}：相域（不平衡）AC/DC 单体式非线性 OPF，入口
        \srcpath{solve\_three\_phase\_hybrid\_opf}
        （\srcpath{include/hacdcpf/optimal\_power\_flow/three\_phase\_hybrid\_opf.hpp:solve\_three\_phase\_hybrid\_opf}），
        其输入 \fld{ThreePhaseHybridOPFCase}（同文件:ThreePhaseHybridOPFCase）是已装配的相域
        算例而非 Rich 系统；从 Rich 系统的转换由适配器
        \srcpath{build\_three\_phase\_hybrid\_model} 完成
        （\srcpath{include/hacdcpf/optimal\_power\_flow/three\_phase\_hybrid\_adapter.hpp:build\_three\_phase\_hybrid\_model}）。
        详见第\ref{sec:threephase}节。
  \item \textbf{AML 建模层}：基于 MIPSolvers AML 的求解器无关建模层
        （\srcpath{include/hacdcpf/power\_models/} 下
        \srcpath{ac\_pf\_model\_builder.hpp}、
        \srcpath{hybrid\_opf\_model\_builder.hpp}、
        \srcpath{dc\_opf\_model\_builder.hpp}、
        \srcpath{lindistflow\_builder.hpp}、
        \srcpath{scuc\_builder.hpp}；实现于 \srcpath{src/power\_models/}），
        供外部求解器与对照实验使用。该实现不属于本模块；其唯一数学规范见
        \srcpath{docs/modules/power\_models/power\_models\_manual.tex}。
  \item \textbf{公共接口层}：选项、问题描述、结果类型与求解器抽象（本章
        第\ref{subsec:overview-options}--\ref{subsec:overview-interface}节）。
\end{itemize}

对外门面 \srcpath{include/hacdcpf/api/hacdcpf.hpp:solve\_ac\_opf} 等四处声明
（:solve\_ac\_opf、:solve\_dc\_opf、:solve\_rpo、:verify\_opf\_result）暴露四个 OPF
相关入口：\allowbreak\srcpath{hacdcpf::solve\_ac\_opf}、\allowbreak\srcpath{hacdcpf::solve\_dc\_opf}、
\allowbreak\srcpath{hacdcpf::solve\_rpo} 与求解后独立校核 \srcpath{verify\_opf\_result}。
门面 \srcpath{solve\_ac\_opf} 在调用内核返回后，额外在 OPF 调度点上重算
开关/断路器端潮流并写入 \fld{ac\_switch\_flows} 与
\fld{ac\_circuit\_breaker\_flows}（\srcpath{src/api/hacdcpf.cpp:solve\_ac\_opf}，
其中支路潮流经 \fld{bus\_merge\_map} 反投影回原始元件，见:1872--1884）；
\srcpath{solve\_dc\_opf} 与 \srcpath{solve\_rpo} 门面是直接转发
（同文件:solve\_dc\_opf）。

\subsubsection*{数据流架构}

\begin{figure}[H]
\centering
\resizebox{\textwidth}{!}{%
\begin{tikzpicture}[
  font=\small,
  box/.style={draw, rectangle, minimum height=0.95cm, align=center, inner sep=4pt},
  arr/.style={-{Stealth[length=2.6mm]}, thick},
  node distance=0.9cm]
  \node[box, minimum width=2.9cm] (rich) {Rich\\ \fld{HybridPowerSystem}};
  \node[box, right=1.0cm of rich, minimum width=2.9cm] (proj) {校验与投影\\ canonical projection};
  \node[box, right=1.0cm of proj, minimum width=2.4cm] (data) {canonical\\ \fld{SolverData}};
  \node[box, right=1.4cm of data, yshift=2.05cm, minimum width=5.1cm] (acopf)
    {AC OPF / 交直流混合 OPF（NLP）\\ Native IPM（纯交流）/ Parity IPM / 嵌入式 Ipopt};
  \node[box, below=0.32cm of acopf, minimum width=5.1cm] (dcopf)
    {DC OPF（线性化 LP/QP）};
  \node[box, below=0.32cm of dcopf, minimum width=5.1cm] (rpo)
    {RPO 无功优化（MINLP）};
  \node[box, below=0.32cm of rpo, minimum width=5.1cm] (tp)
    {三相混合 OPF（相域）};
  \node[box, below=0.32cm of tp, minimum width=5.1cm] (aml)
    {AML 建模层（\srcpath{src/power\_models/}）};
  \node[box, right=1.4cm of dcopf, minimum width=3.6cm] (res)
    {结果归因（反投影）\\ \fld{ACOPFResult} / \fld{DCOPFResult}};
  \node[box, below=0.55cm of res, minimum width=3.6cm] (audit)
    {独立校核\\ \srcpath{verify\_opf\_result}};
  \draw[arr] (rich) -- (proj);
  \draw[arr] (proj) -- (data);
  \draw[arr] (data.east) -- (acopf.west);
  \draw[arr] (data.east) -- (dcopf.west);
  \draw[arr] (data.east) -- (rpo.west);
  \draw[arr] (rich.south) |- (tp.west);
  \draw[arr] (rich.south) |- (aml.west);
  \draw[arr] (acopf.east) -- (res.north west);
  \draw[arr] (dcopf.east) -- (res.west);
  \draw[arr] (rpo.east) -- (res.south west);
  \draw[arr] (res) -- (audit);
\end{tikzpicture}}
\caption{OPF 求解器家族数据流。Rich 模型经校验与 canonical 投影装配为
\fld{SolverData} 后送入各求解引擎；结果经映射归因回原始组件稳定 ID，并可由
\srcpath{verify\_opf\_result} 做独立于求解器的可行性校核。三相混合 OPF 与
AML 建模层不经 \fld{SolverData} 主通道，前者经适配器直接装配相域算例，后者
自建 AML 模型。}
\label{fig:overview-family}
\end{figure}

数据流遵循全平台统一主线：Rich \fld{HybridPowerSystem}（工程语义层）
$\rightarrow$ 校验与 canonical 投影（开关收缩、零阻抗合并，产生
\fld{BusMergeMap}/\fld{BranchExpandMap}）$\rightarrow$ 装配 \fld{SolverData}
$\rightarrow$ 求解 $\rightarrow$ 结果反投影归因回原始组件稳定 ID。DC OPF
显式复用同一投影算子
（\srcpath{src/optimal\_power\_flow/dc\_opf.cpp:solve\_dc\_opf}，调用
\srcpath{projection::RichToCanonicalOperator::apply}），其母线/支路结果经
\srcpath{unproject\_bus\_vector} 与 \fld{branch\_orig\_to\_proj} 映射回原
始 ID 空间（同文件:separate\_external\_grid\_dispatch（lambda））；AC OPF 的结果归因集中在
\srcpath{separate\_external\_grid\_dispatch}
（\srcpath{src/optimal\_power\_flow/ac\_opf.cpp:separate\_external\_grid\_dispatch（lambda）}），外部电网注入被
从扩展发电机数组中剥离回填 \fld{external\_grid\_p\_mw}/
\fld{external\_grid\_q\_mvar}（:separate\_external\_grid\_dispatch（lambda））。

需要如实说明的三点：

\begin{itemize}
  \item \textbf{DC OPF 拒绝 LCC 算例}：只要系统中存在在运 LCC 换流站，
        \srcpath{solve\_dc\_opf} 直接返回拒绝状态而不求解，\fld{status} 为
        ``DC OPF rejected: this AC-only linear formulation does not model
        LCC AC/DC coupling\ldots''，\fld{model\_scope} 置
        \fld{dc-opf:rejected-lcc}
        （\srcpath{src/optimal\_power\_flow/dc\_opf.cpp:solve\_dc\_opf}）。
  \item \textbf{混合算例被强制走向全空间路径}：只要系统含 DC/VSC 子系统
        或 LCC，AC OPF 在路径分派处强制打开 \fld{enable\_primal\_dual} 与
        \fld{use\_parity\_ipm}（除非显式选择了 \fld{EconomicDispatch} 后端），
        即传统 Native AC 路径实际上只服务纯交流算例
        （\srcpath{src/optimal\_power\_flow/ac\_opf.cpp:solve\_ac\_opf\_impl}；混合标志
        计算于:3261--3265）。显式选择 \fld{EconomicDispatch} 的混合算例
        不会被静默降级，而是直接报错并给出提示
        （同文件:solve\_ac\_opf\_impl）。
  \item \textbf{AML 建模层与 Parity 建模层是两套独立装配}：Parity IPM 与
        嵌入式 Ipopt 共用 \srcpath{parity::build\_problem} 装配的全空间
        模型（\srcpath{src/optimal\_power\_flow/ac\_opf.cpp:solve\_with\_parity\_ipm}），并不
        经过 \srcpath{src/power\_models/} 的 AML builder。
        （\srcpath{src/optimal\_power\_flow/ac\_opf.cpp:solve\_with\_parity\_ipm} 的注释
        称“Parity-IPM uses the AML hybrid OPF builder”，与:2540 的实际
        调用不符，本章以代码为准。）
\end{itemize}

\subsection{非线性引擎对比：Native AC 与交直流混合两条路线}\label{subsec:overview-engines}

\fld{ACOPFOptions::ac\_solver\_backend} 在四个取值中选择非线性引擎
（枚举定义于
\srcpath{include/hacdcpf/optimal\_power\_flow/opf\_options.hpp:ACOPFSolverBackend}）：
\fld{Auto}（默认，Parity 全空间 IPM 优先、不收敛且编译了 Ipopt 时回退
Ipopt）、\fld{ParityIPM}（自研全空间原始对偶 IPM）、\fld{Ipopt}（嵌入式
Ipopt，滤波线搜索 NLP 求解器，作用于同一 Parity 建模）、
\fld{EconomicDispatch}（merit-order 经济调度加一次 AC 潮流，快速但次优、
仅限纯交流系统）。分派逻辑把该枚举翻译成 legacy 标志
\fld{enable\_primal\_dual}/\fld{use\_parity\_ipm} 与内部枚举
\srcpath{ParityInnerSolver}
（\srcpath{src/optimal\_power\_flow/ac\_opf.cpp:solve\_ac\_opf\_impl}；内部枚举定义于
:1870--1873）。其中 Parity IPM 与嵌入式 Ipopt 作用于同一份 Parity 全空间
建模，是\textbf{交直流混合 OPF 的两条承载路线}；Native AC 紧致路径只承担
纯交流算例。三条非线性引擎的对比见表~\ref{tab:overview-engines}。

\begin{table}[H]
\centering
\footnotesize
\begin{tabular}{@{}p{2.1cm}p{3.6cm}p{3.9cm}p{3.9cm}@{}}
\toprule
\textbf{维度} & \textbf{Native AC（传统路径）} & \textbf{Parity IPM（自研全空间）} & \textbf{嵌入式 Ipopt}\\
\midrule
建模口径 &
降阶传统建模：\srcpath{core::SolverData} + \fld{ACOPFIndex}，潮流雅可比上下文
（\srcpath{src/optimal\_power\_flow/ac\_opf.cpp:solve\_ac\_opf\_impl}） &
Parity 全空间建模：全部电压/功率/换流器变量为显式原始变量
（\srcpath{include/hacdcpf/optimal\_power\_flow/formulation.hpp:Problem}） &
与 Parity IPM 完全同一模型，1:1 映射为 \fld{engine::NLPModel}
（\srcpath{src/optimal\_power\_flow/ac\_opf.cpp:build\_parity\_nlp\_input}）\\
\midrule
坐标与变量布局 &
极坐标 $(\theta,V)$ 加发电机功率；非线性工程限值经外部栅栏/罚外环处理
（结果 \fld{model\_limitations} 自述，同文件:separate\_external\_grid\_dispatch（lambda）） &
极坐标 AC + 直流电压 + 换流器/储能/可控资源全变量块，布局见 \fld{VarIndex}
（\srcpath{include/hacdcpf/optimal\_power\_flow/formulation.hpp:VarIndex}） &
同 Parity IPM（共享 \fld{Problem::vidx}/\fld{cidx}）\\
\midrule
求解核 &
内环 KKT 牛顿 + 外环罚迭代；\fld{solver\_path=NativeAC}
（\srcpath{src/optimal\_power\_flow/ac\_opf.cpp:solve\_ac\_opf\_impl}） &
自研原始对偶 IPM \srcpath{solve\_primal\_dual\_ipm}
（\srcpath{include/hacdcpf/optimal\_power\_flow/native\_ipm\_solver.hpp:solve\_primal\_dual\_ipm}）；
\fld{solver\_path=ParityIPM}（src/optimal\_power\_flow/ac\_opf.cpp:solve\_with\_parity\_ipm） &
MIPSolvers 内嵌 Ipopt（滤波线搜索），经 \fld{engine::IpoptAdapter}
（\srcpath{src/optimal\_power\_flow/ac\_opf.cpp:solve\_parity\_with\_ipopt}）\\
\midrule
换流器/LCC 建模 &
换流器为自由 $P_{ac}/Q_{ac}/P_{dc}$ 盒式变量 + 容量圆，不钉 VDC\_Q 直流
电压（\fld{model\_scope}=\allowbreak\texttt{ac-opf-\allowbreak native:\allowbreak vsc-free-\allowbreak pq+\allowbreak capacity-\allowbreak circle}，
同文件:solve\_ac\_opf\_impl）；不支持 LCC（被强制转走，:solve\_ac\_opf\_impl） &
VDC\_Q 直流电压等式恒强制，容量圆/电流/调制度约束族可开关
（\fld{model\_scope} 标签按实际强制的约束族拼接，同文件:solve\_with\_parity\_ipm）；
支持 LCC 准稳态（:solve\_with\_parity\_ipm） &
与 Parity IPM 相同（共享建模）\\
\midrule
节点电价 LMP &
不提取节点电价：\fld{lmp\_p}/\fld{lmp\_q} 留空、\fld{lmp\_valid=false}，
由结果归因层统一补写原因（同文件:separate\_external\_grid\_dispatch（lambda）） &
从等式对偶 $\lambda$ 提取 \fld{lmp\_p}/\fld{lmp\_q} 并置 \fld{lmp\_valid}
（同文件:solve\_with\_parity\_ipm） &
适配器不回传约束乘子，\fld{lmp\_valid=false} 并注明原因
（同文件:solve\_with\_parity\_ipm）\\
\midrule
编译门控与可用性 &
始终可用 &
始终可用（自研代码） &
受 \srcpath{HACDCPF\_HAVE\_IPOPT} 门控；未编译时 \fld{available=false}、
\texttt{status=``Ipopt not compiled in''}
（同文件:solve\_parity\_with\_ipopt），显式请求时回退 Parity IPM 并在后端标签注明
（:solve\_with\_parity\_ipm）。详见第\ref{subsec:overview-ipopt}节\\
\midrule
适用场景 &
纯交流中小算例、回归对照；混合算例不可用 &
生产默认：交直流混合/LCC 算例的统一路径 &
大系统/困难算例的 Auto 回退与独立对照\\
\bottomrule
\end{tabular}
\caption{三条 AC 非线性引擎对比。第四条路径 \fld{EconomicDispatch} 为
merit-order 经济调度加一次牛顿 AC 潮流（实现
\srcpath{src/optimal\_power\_flow/ac\_opf.cpp:try\_dispatch\_pf\_fallback}，入口
:3325--3329、3520--3525），不是非线性优化，故不列入。}
\label{tab:overview-engines}
\end{table}

Auto 模式的完整分派顺序如下（
\srcpath{src/optimal\_power\_flow/ac\_opf.cpp}）：
\begin{enumerate}
  \item \fld{ac\_solver\_backend} 枚举翻译成 legacy 标志与
        \srcpath{ParityInnerSolver}（:solve\_ac\_opf\_impl）；
  \item 含 DC/VSC 或 LCC 的算例强制启用 Parity 全空间路径（:solve\_ac\_opf\_impl）；
  \item \fld{objective\_homotopy} 为真且非 Ipopt 后端时走两阶段目标同伦
        （$t:0\to1$ 的成本同伦 + Davidenko 切向预测，:solve\_ac\_opf\_impl）；显式
        Ipopt 与同伦不组合，退化为独立单解（:solve\_ac\_opf\_impl）；
  \item Parity 路径内部再按 \srcpath{ParityInnerSolver} 选内层牛顿引擎：
        \fld{Auto} 且启用 \fld{ac\_pf\_warm\_start} 时先用独立 Ipopt 再以
        原生 IPM 精化（:solve\_with\_parity\_ipm）；显式 \fld{Ipopt} 但不可用时回退原生
        IPM（:solve\_with\_parity\_ipm）；其余 \fld{Auto} 情形先原生 Parity IPM，不收敛
        且 \fld{allow\_fallback} 时回退 Ipopt（:solve\_with\_parity\_ipm）；
  \item 未启用原始对偶路径时：混合算例报错（:solve\_ac\_opf\_impl），纯交流算例走
        经济调度快速路径（:solve\_ac\_opf\_impl）或传统 Native AC 路径
        （:solve\_ac\_opf\_impl）。
\end{enumerate}

\subsection{公共选项：ACOPFOptions 与 DCOPFOptions}\label{subsec:overview-options}

\compmeta{ACOPFOptions}{\srcpath{include/hacdcpf/optimal\_power\_flow/opf\_options.hpp:ACOPFOptions}}

\fld{ACOPFOptions} 是 AC OPF（三条引擎共用）的全部用户选项。注意两点实现
口径：其一，\srcpath{solve\_ac\_opf} 入口处会把非正（越界）的数值字段重置
为默认值（\srcpath{src/optimal\_power\_flow/ac\_opf.cpp:solve\_ac\_opf\_impl}），因此
传 0 等价于取默认；其二，\fld{ac\_solver\_backend} 与 legacy 标志
\fld{enable\_primal\_dual}/\fld{use\_parity\_ipm} 并存，\fld{Auto} 时由
legacy 标志继续控制路径选择（字段注释，
\srcpath{include/hacdcpf/optimal\_power\_flow/opf\_options.hpp:ac\_solver\_backend}）。

\begin{paramtable}
\fld{ac\_solver\_backend} & — & — & Auto/ParityIPM/\allowbreak Ipopt/\allowbreak EconomicDispatch & Auto & 非线性引擎选择（枚举 \fld{ACOPFSolverBackend}，include/hacdcpf/optimal\_power\_flow/opf\_options.hpp:ACOPFSolverBackend）\\
\fld{objective} & — & — & Economic/\allowbreak VoltageDeviation/\allowbreak ActiveLoss/\allowbreak VoltageDeviation\allowbreak AndLoss & Economic & 连续目标（枚举 \fld{ACOPFObjective}，:ACOPFObjective）；后三者供 RPO 内层 NLP 使用\\
\fld{voltage\_target\_pu} & $V_{\mathrm{ref}}$ & pu & 0.95--1.05（经验值） & 1.0 & 电压偏差目标值\\
\fld{voltage\_deviation\_weight} & $w_V$ & — & $\ge 0$ & 1.0 & 电压偏差目标项权重\\
\fld{active\_loss\_weight} & $w_{\mathrm{loss}}$ & — & $\ge 0$ & 1.0 & 有功网损目标项权重\\
\fld{max\_inner\_iterations} & — & — & 1--$10^{5}$（HTTP 层 clamp 至 1--100000，\srcpath{runtime\_api\_v1.cpp:ac\_opf\_options}） & 80 & IPM 内层最大迭代数\\
\fld{max\_outer\_iterations} & — & — & 1--$10^{4}$（HTTP 层 clamp，\srcpath{runtime\_api\_v1.cpp:ac\_opf\_options}） & 8 & 外环（罚/栅栏）最大迭代数\\
\fld{max\_line\_search\_steps} & — & — & $\ge 1$ & 20 & 单次线搜索最大步数\\
\fld{feasibility\_tol} & $\varepsilon_{\mathrm{feas}}$ & pu & $10^{-12}$--1（HTTP 层 clamp，\srcpath{runtime\_api\_v1.cpp:ac\_opf\_options}） & $10^{-6}$ & 原始可行性容差\\
\fld{stationarity\_tol} & $\varepsilon_{\mathrm{stat}}$ & pu & 同上 & $10^{-6}$ & 驻点/对偶容差\\
\fld{barrier\_mu0} & $\mu_0$ & — & $>0$（典型 $10^{-3}$--$10^{-1}$） & $10^{-2}$ & 初始栅栏参数\\
\fld{barrier\_mu\_reduction} & — & — & $(0,1)$ & 0.2 & 栅栏参数缩减因子\\
\fld{barrier\_mu\_min} & $\mu_{\min}$ & — & $>0$ & $10^{-8}$ & 栅栏参数下限\\
\fld{merit\_penalty} & — & — & $>0$ & 10.0 & merit 函数罚系数（传统路径）\\
\fld{regularization} & — & — & $>0$ & $10^{-6}$ & KKT 系统正则化\\
\fld{step\_backoff} & — & — & $(0,1)$ & 0.5 & 线搜索步长回退因子\\
\fld{interior\_fraction} & — & — & $(0,1)$ & 0.995 & fraction-to-boundary 内点保持系数\\
\fld{ac\_eval\_threads} & — & — & $\ge 1$ & 1 & AC 方程/雅可比求值线程数\\
\fld{enable\_primal\_dual} & — & — & true/false & false & legacy 标志：启用原始对偶路径\\
\fld{use\_parity\_ipm} & — & — & true/false & false & legacy 标志：选用 Parity 全空间建模\\
\fld{allow\_fallback} & — & — & true/false & true & 允许 Auto 模式在不收敛时回退（Ipopt 或同伦末段全目标重试）\\
\fld{verbose} & — & — & true/false & false & 迭代日志\\
\fld{enable\_phase\_one} & — & — & true/false & true & Parity Native IPM 的有界 primal/dual warm-start 构造\\
\fld{phase\_one\_time\_limit\_ms} & — & ms & $>0$ 协作截止；$\le0$ 不限 & 5000 & Phase I 墙钟预算\\
\fld{phase\_one\_max\_iterations} & — & — & $\ge0$ & 12 & Phase I Newton 迭代硬上限\\
\fld{phase\_one\_max\_factorizations} & — & — & $\ge0$ & 14 & Phase I primal/dual 共用分解上限\\
\fld{phase\_one\_max\_backtracks} & — & — & $\ge0$ & 12 & 每次 Phase I 步的回溯上限\\
\fld{phase\_one\_barrier\_mu} / \fld{phase\_one\_admission\_mu\_factor} & $\mu_0,\eta_a$ & — & $(0,0.1]$ / $\ge0$ & 0.1 / 1.0 & central warm-start 的障碍尺度与原坐标准入比例\\
\fld{phase\_one\_primal\_mu\_factor} / \fld{phase\_one\_centrality\_tolerance} & $\eta_p,\eta_c$ & — & $\ge0$ & 0.1 / 0.5 & primal handoff corridor 与 $\max_i|s_i z_i/\mu_0-1|$ 门限\\
\fld{warm\_start} & — & — & 指针/nullptr & nullptr & 指向前次 \fld{ACOPFResult} 的非拥有暖启动指针；维度不兼容的变量族各自回退物理初值（字段注释，include/hacdcpf/optimal\_power\_flow/opf\_options.hpp:ACOPFOptions）\\
\fld{ac\_pf\_warm\_start} & — & — & true/false & false & 先解一次混合潮流，以 AC 可行点初始化 Parity IPM（字段注释：PEGASE 受压算例由此解锁，:ac\_pf\_warm\_start；实现 src/optimal\_power\_flow/ac\_opf.cpp:solve\_with\_parity\_ipm）\\
\fld{ac\_pf\_dc\_phase\_one} & — & — & true/false & true & 已启用 AC-PF warm start 时，对至少 5000 母线的纯 AC 模型按 baseline-dual 预筛选运行有限步 compact NativeLCQP dispatch seed\\
\fld{ac\_pf\_dc\_phase\_one\_time\_limit\_ms} & $T_{\mathrm{DC},I}$ & ms & $>0$ 协作截止；$\le0$ 不限 & 2000 & DC Phase I 端到端预算，含 projection/formulation/structural/symbolic/numeric；允许一个在途稀疏分解超时\\
\fld{ac\_pf\_dc\_phase\_one\_max\_iterations} / \fld{ac\_pf\_dc\_phase\_one\_tolerance} & — & — / pu & $\ge1$ / $\ge0$ & 5 / $10^{-2}$ & DC warm-start-only 迭代预算与可用残差门限\\
\fld{ac\_pf\_dc\_phase\_one\_min\_dual\_improvement} / \fld{ac\_pf\_dc\_phase\_one\_baseline\_dual\_threshold} & — & — & $[0,1]$ / $\ge0$ & 0.2 / 100 & 候选 dual 最小相对改善与是否支付 DC Phase I 的 baseline 门限\\
\fld{compute\_homotopy\_tangent} & — & — & true/false & false & 在返回点计算 Davidenko 目标同伦切向量 $\mathrm{d}w/\mathrm{d}t$（一次额外惯性受控 KKT 解，:compute\_homotopy\_tangent）\\
\fld{homotopy\_t} & $t$ & — & $(0,1]$ & 1.0 & 所解问题的成本尺度 $t$（切向量右端为 $\nabla f/t$）\\
\fld{objective\_homotopy} & — & — & true/false & false & 两阶段目标同伦求解（$t:0\to1$，:objective\_homotopy；实现 src/optimal\_power\_flow/ac\_opf.cpp:solve\_ac\_opf\_impl）\\
\fld{homotopy\_dt0} & $\Delta t_0$ & — & $(0,1]$ & 0.10 & 初始延拓步长；非正值回退 0.10（src/optimal\_power\_flow/ac\_opf.cpp:solve\_ac\_opf\_impl）\\
\fld{enforce\_branch\_limits} & — & — & true/false & true & 约束族开关：支路容量\\
\fld{enforce\_converter\_capacity} & — & — & true/false & true & 约束族开关：换流器容量圆\\
\fld{enforce\_converter\_current\_limits} & — & — & true/false & true & 约束族开关：换流器电流限值\\
\fld{enforce\_converter\_modulation\_limits} & — & — & true/false & true & 约束族开关：调制度窗口与 DC/DC 占空比\\
\end{paramtable}

\compmeta{DCOPFOptions}{\srcpath{include/hacdcpf/optimal\_power\_flow/opf\_options.hpp:DCOPFOptions}}

\fld{DCOPFOptions} 是 DC OPF 的全部用户选项。后端枚举
\fld{DCOPFSolverBackend}（:DCOPFSolverBackend）取
\fld{Auto}/\fld{Native}/\fld{NativeQP}/\fld{HiGHS}/\fld{Gurobi}；实际回退
链见第\ref{sec:dcopf}节，Auto 链为 Gurobi-QP $\to$ NativeLCQP $\to$ HiGHS
$\to$ NativeDualSimplex
（\srcpath{src/optimal\_power\_flow/dc\_opf.cpp:solve\_dc\_opf}），每步尝试记录
进 \fld{solver\_chain}（:solve\_dc\_opf）。

\begin{paramtable}
\fld{feasibility\_tol} & $\varepsilon_{\mathrm{feas}}$ & pu & $10^{-12}$--1（HTTP 层 clamp，\srcpath{runtime\_api\_v1.cpp:dc\_opf\_options}） & $10^{-6}$ & LP/QP 可行容差\\
\fld{optimality\_tol} & $\varepsilon_{\mathrm{opt}}$ & pu & $>0$ & $10^{-6}$ & 最优性（对偶）容差\\
\fld{max\_iterations} & — & — & 1--$10^{6}$（HTTP 层 clamp，\srcpath{runtime\_api\_v1.cpp:dc\_opf\_options}） & 10000 & 单纯形/IPM 最大迭代数\\
\fld{verbose} & — & — & true/false & false & 日志\\
\fld{solver} & — & — & Auto/Native/\allowbreak NativeQP/\allowbreak HiGHS/\allowbreak Gurobi & Auto & LP/QP 后端选择（枚举 \fld{DCOPFSolverBackend}，include/hacdcpf/optimal\_power\_flow/opf\_options.hpp:DCOPFSolverBackend）\\
\fld{pwl\_segments} & — & — & 1--1000（HTTP 层 clamp，\srcpath{runtime\_api\_v1.cpp:dc\_opf\_options}） & 4 & 二次成本分段线性化段数（LP 路径）\\
\fld{include\_branch\_limits} & — & — & true/false & true & 是否计入支路潮流限值\\
\fld{branch\_limit\_margin} & — & pu & $>0$ & 1.0 & 支路限值缩放系数\\
\fld{load\_shedding} & — & — & true/false & true & 是否启用切负荷松弛变量\\
\fld{voll} & VOLL & — & $\ge 0$ & 0.0 & 失负荷价值（成本币种未定）；$\le 0$ 时按 $\max(100,\,1.1\times$ 最大边际成本$)$ 自动折算（\srcpath{dc\_opf.cpp:build\_dc\_opf\_lp}）\\
\fld{compute\_lmp} & — & — & true/false & true & 是否在主解后补解支撑 LP 以恢复约束对偶（LMP 与阻塞价格）；关闭可省一次 LP（字段注释，include/hacdcpf/optimal\_power\_flow/opf\_options.hpp:DCOPFOptions；实现 src/optimal\_power\_flow/dc\_opf.cpp:populate\_missing\_duals\_from\_supporting\_lp）\\
\fld{structural\_warm\_start} & — & — & true/false & true & NativeLCQP 的逐连通分量配平与约化 Laplacian 初值；其他后端忽略\\
\end{paramtable}

\subsection{结果类型与诚实口径}\label{subsec:overview-result}

\subsubsection*{ACOPFResult 字段}

\compmeta{ACOPFResult}{\srcpath{include/hacdcpf/optimal\_power\_flow/opf\_result.hpp:ACOPFResult}}

结果向量的索引空间一律为\textbf{原始（authored）组件顺序}：母线向量按
\fld{sys.ac.buses}/\fld{sys.dc.buses} 顺序，设备向量按对应容器顺序；归因
映射经 \texttt{ComponentRef\{original\_index, source\_type\}}（:ComponentRef）
回指组件稳定 \fld{index}。量纲除注明外为有名值（MW/MVAr/kV 制下的 pu）。

\begin{paramtable}
\fld{vm} & $V_m$ & pu & 约 0.9--1.1 & 空 & AC 母线电压幅值（authored bus 顺序）\\
\fld{va} & $\theta$ & rad & $-\pi$--$\pi$ & 空 & AC 母线电压相角（HTTP 序列化名 \fld{va\_rad}，\srcpath{runtime\_api\_v1.cpp:serialize\_ac\_opf}）\\
\fld{vdc} & $V_{dc}$ & pu & 约 0.9--1.1 & 空 & DC 母线电压（标幺）\\
\fld{pg\_mw} & $P_g$ & MW & $P_{\min}$--$P_{\max}$ & 空 & 发电机有功（原始发电机顺序，外部电网已剥离）\\
\fld{qg\_mvar} & $Q_g$ & MVAr & $Q_{\min}$--$Q_{\max}$ & 空 & 发电机无功\\
\fld{external\_grid\_p\_mw} & $P_{\mathrm{eg}}$ & MW & — & 空 & 外部电网有功（由扩展发电机行剥离回填，src/optimal\_power\_flow/ac\_opf.cpp:separate\_external\_grid\_dispatch（lambda））\\
\fld{external\_grid\_q\_mvar} & $Q_{\mathrm{eg}}$ & MVAr & — & 空 & 外部电网无功\\
\fld{dpd\_mw} & $\Delta P^d$ & MW & $\ge 0$ & 空 & 各母线有功切负荷量\\
\fld{dqd\_mvar} & $\Delta Q^d$ & MVAr & $\ge 0$ & 空 & 各母线无功切负荷量\\
\fld{pac\_mw} & $P^{ac}_c$ & MW & $\pm S_N$ & 空 & VSC 交流侧有功\\
\fld{qac\_mvar} & $Q^{ac}_c$ & MVAr & $\pm S_N$ & 空 & VSC 交流侧无功\\
\fld{lcc\_transfers} & — & — & — & 空 & LCC 准稳态运行点（按 \fld{LCCConverter::index} 键控，字段注释 include/hacdcpf/optimal\_power\_flow/opf\_result.hpp:ACOPFResult）\\
\fld{pren\_mw} / \fld{qren\_mvar} & $P^{ren},Q^{ren}$ & MW/MVAr & $\ge 0$ & 空 & 可再生能源出力\\
\fld{pstor\_mw} / \fld{qstor\_mvar} & $P^{stor},Q^{stor}$ & MW/MVAr & $\pm S_N$ & 空 & 储能充放功率\\
\fld{pdcdc\_mw} & $P^{dcdc}$ & MW & $\pm P_N$ & 空 & DC/DC 变换器传输功率\\
\fld{pflex\_mw} & $P^{flex}$ & MW & — & 空 & 柔性负荷调节量\\
\fld{er\_port\_p\_mw} / \fld{er\_port\_q\_mvar} & $P^{er},Q^{er}$ & MW/MVAr & — & 空 & 能量路由器端口功率\\
\fld{lmp\_p} & $\lambda^P$ & — & — & 空 & AC 有功节点电价（币种未定；由等式对偶换算，src/optimal\_power\_flow/ac\_opf.cpp:solve\_with\_parity\_ipm）\\
\fld{lmp\_q} & $\lambda^Q$ & — & — & 空 & AC 无功节点电价\\
\fld{lmp\_valid} & — & — & true/false & false & 仅当价格来自本解所出建模的认证等式乘子时为真（字段注释，include/hacdcpf/optimal\_power\_flow/opf\_result.hpp:ACOPFResult）\\
\fld{lmp\_validity\_reason} & — & — & — & 空 & \fld{lmp\_valid} 判定的人读原因\\
\fld{ipm\_primal\_state} & $x$ & — & — & 空 & 原生 IPM 延拓状态（原始变量），供 \fld{warm\_start} 复用（字段注释 :109--112）\\
\fld{ipm\_equality\_dual\_state} & $\lambda$ & — & — & 空 & 延拓状态：等式对偶\\
\fld{ipm\_inequality\_dual\_state} & $\mu$ & — & — & 空 & 延拓状态：不等式对偶\\
\fld{ipm\_slack\_state} & $s$ & — & — & 空 & 延拓状态：松弛变量\\
\fld{ipm\_tangent\_primal} & $\mathrm{d}x/\mathrm{d}t$ & — & — & 空 & Davidenko 切向量（原始变量块），仅 \fld{compute\_homotopy\_tangent} 时填充（:ACOPFResult）\\
\fld{ipm\_tangent\_slack} & $\mathrm{d}s/\mathrm{d}t$ & — & — & 空 & 切向量：松弛块\\
\fld{ipm\_tangent\_equality\_dual} & $\mathrm{d}\lambda/\mathrm{d}t$ & — & — & 空 & 切向量：等式对偶块\\
\fld{ipm\_tangent\_inequality\_dual} & $\mathrm{d}\mu/\mathrm{d}t$ & — & — & 空 & 切向量：不等式对偶块\\
\fld{gen\_map} & — & — & — & 空 & 发电机归因映射（\fld{original\_index}+\fld{source\_type}）\\
\fld{ren\_map} & — & — & — & 空 & 可再生能源归因映射\\
\fld{stor\_map} & — & — & — & 空 & 储能归因映射\\
\fld{dcdc\_map} & — & — & — & 空 & DC/DC 归因映射\\
\fld{flex\_map} & — & — & — & 空 & 柔性负荷归因映射\\
\fld{er\_port\_map} & — & — & — & 空 & 能量路由器端口归因映射\\
\fld{converged} & — & — & true/false & false & 收敛标记（Ipopt 路径含“可接受残差”判据，src/optimal\_power\_flow/ac\_opf.cpp:solve\_parity\_with\_ipopt）\\
\fld{iterations} & — & — & $\ge 0$ & 0 & 总迭代数\\
\fld{outer\_iterations} & — & — & $\ge 0$ & 0 & 外环迭代数（传统路径）\\
\fld{objective} & $f$ & — & $\ge 0$ & 0 & 目标值（成本币种未定）\\
\fld{max\_constraint\_violation} & — & pu & $\ge 0$ & 0 & 返回点最大约束违反\\
\fld{max\_stationarity} & — & pu & $\ge 0$ & 0 & 返回点驻点残差\\
\fld{status} & — & — & — & 空 & 人读状态串（引擎标签原样保留，如 ``Ipopt: \ldots''）\\
\fld{solver\_path} & — & — & Unknown/\allowbreak NativeAC/\allowbreak ParityIPM & Unknown & 求解路径枚举（include/hacdcpf/optimal\_power\_flow/opf\_result.hpp:OPFSolverPath）。注意 Ipopt 解也记 \fld{ParityIPM}，真实内层引擎见 \fld{profiling.linear\_solver\_backend}（src/optimal\_power\_flow/ac\_opf.cpp:solve\_with\_parity\_ipm）；经济调度快速路径不置该字段，保持 \fld{Unknown}\\
\fld{profiling} & — & — & — & 默认 & 求解剖析；含 KKT 计数/残差及 Phase I 前后 violation、dual-fit、预算、终止、线性后端与 Phase II 接受状态\\
\fld{infeasibility\_hints} & — & — & — & 空 & 不可行/不收敛的可能原因（\fld{converged=false} 时由求解器和/或 \srcpath{verify\_opf\_result} 填充，include/hacdcpf/optimal\_power\_flow/opf\_result.hpp:ACOPFResult）\\
\fld{audit} & — & — & — & 默认 & 独立求解后校核（\fld{OpfAudit}，见下表）\\
\fld{converter\_model\_scope} & — & — & — & 默认 & 换流器模型保真度声明（见下“诚实口径”）\\
\fld{model\_limitations} & — & — & — & 空 & 本求解路径的能力边界声明（见下“诚实口径”）\\
\fld{ac\_switch\_flows} & — & — & — & 空 & OPF 调度点上的开关端潮流（门面层填充，src/api/hacdcpf.cpp:solve\_ac\_opf）\\
\fld{ac\_circuit\_breaker\_flows} & — & — & — & 空 & OPF 调度点上的断路器端潮流\\
\end{paramtable}

\subsubsection*{OpfAudit 字段}

\compmeta{OpfAudit}{\srcpath{include/hacdcpf/optimal\_power\_flow/opf\_result.hpp:OpfAudit}}

\fld{OpfAudit} 是独立于求解器的求解后可行性审计，由门面函数
\srcpath{verify\_opf\_result} 填充（声明
\srcpath{include/hacdcpf/api/hacdcpf.hpp:verify\_opf\_result}；实现
\srcpath{src/api/hacdcpf.cpp:verify\_opf\_result}）。HTTP v1 接口在序列化 AC OPF
结果前自动执行该审计（\srcpath{src/server/runtime\_api\_v1.cpp:serialize\_ac\_opf}）。

\begin{paramtable}
\fld{audited} & — & — & true/false & false & 审计已执行标记（src/api/hacdcpf.cpp:verify\_opf\_result）\\
\fld{max\_power\_balance\_violation\_mw} & — & MW & $\ge 0$ & 0 & 母线功率平衡最大违反。\textbf{当前实现未填充}：\srcpath{verify\_opf\_result} 只检查电压限值、发电机限值与目标一致性（src/api/hacdcpf.cpp:verify\_opf\_result），该字段保持默认 0\\
\fld{max\_voltage\_limit\_violation\_pu} & — & pu & $\ge 0$ & 0 & 电压越限最大值（src/api/hacdcpf.cpp:verify\_opf\_result）\\
\fld{max\_branch\_limit\_violation\_pu} & — & pu & $\ge 0$ & 0 & 支路负载越限最大值。\textbf{当前实现未填充}（同上）\\
\fld{max\_gen\_limit\_violation\_mw} & — & MW & $\ge 0$ & 0 & 发电机出力越限最大值（src/api/hacdcpf.cpp:verify\_opf\_result）\\
\fld{objective\_recomputed} & — & — & $\ge 0$ & 0 & 由 \fld{pg\_mw}/\fld{qg\_mvar} 重算的目标（src/api/hacdcpf.cpp:verify\_opf\_result）\\
\fld{objective\_reported} & — & — & $\ge 0$ & 0 & 求解器上报的目标（src/api/hacdcpf.cpp:verify\_opf\_result）\\
\fld{objective\_discrepancy\_pct} & — & \% & $\ge 0$ & 0 & 重算与上报目标的相对偏差（src/api/hacdcpf.cpp:verify\_opf\_result）\\
\fld{violations} & — & — & — & 空 & 人读违约描述列表\\
\fld{feasible()} & — & — & true/false & — & 成员函数：\fld{violations} 为空即真（include/hacdcpf/optimal\_power\_flow/opf\_result.hpp:OpfAudit）\\
\end{paramtable}

需要如实指出：头文件注释（include/hacdcpf/api/hacdcpf.hpp:verify\_opf\_result）称该审计“Checks power
balance, voltage limits, branch limits, generator limits, and objective
consistency”，但当前实现仅覆盖电压限值、发电机限值与目标一致性三项
（另在 \fld{converged=false} 时追加无源负荷、限值冲突等
\fld{infeasibility\_hints}，src/api/hacdcpf.cpp:verify\_opf\_result）；功率平衡与支路限值
两个审计字段当前恒为 0，消费方不得据此断言对应可行性。

\subsubsection*{DCOPFResult 字段}

\compmeta{DCOPFResult}{\srcpath{include/hacdcpf/optimal\_power\_flow/opf\_result.hpp:DCOPFResult}}

\begin{paramtable}
\fld{va} & $\theta$ & rad & $-\pi$--$\pi$ & 空 & AC 母线相角（authored bus 顺序）\\
\fld{pg\_mw} & $P_g$ & MW & $P_{\min}$--$P_{\max}$ & 空 & 发电机有功\\
\fld{external\_grid\_p\_mw} & $P_{\mathrm{eg}}$ & MW & — & 空 & 外部电网有功（提升为 slack 发电机后剥离回填，src/optimal\_power\_flow/dc\_opf.cpp:separate\_external\_grid\_dispatch（lambda））\\
\fld{pf\_mw} & $P_f$ & MW & $\pm S_{\max}$ & 空 & 支路有功潮流（映射回原始支路顺序，src/optimal\_power\_flow/dc\_opf.cpp:separate\_external\_grid\_dispatch（lambda））\\
\fld{converged} & — & — & true/false & false & 收敛标记\\
\fld{iterations} & — & — & $\ge 0$ & 0 & 迭代数\\
\fld{objective} & $f$ & — & $\ge 0$ & 0 & 目标值（语义见 \fld{objective\_model}）\\
\fld{status} & — & — & — & 空 & 人读状态串\\
\fld{solver\_name} & — & — & — & 空 & 实际求解后端名（如 \fld{NativeLCQP}、\fld{HiGHS-LP(QP-fallback)}，src/optimal\_power\_flow/dc\_opf.cpp:extract\_dc\_opf\_result、1105）\\
\fld{runtime\_sec} & — & s & $\ge 0$ & 0 & 墙钟耗时\\
\fld{structural\_warm\_start\_*} / \fld{solver\_initial\_primal\_residual} & — & pu/计数/布尔 & — & inf/0/false & NativeLCQP 结构初值的请求、构建、消费、分量/分解、原坐标 residual 和第 0 次迭代 residual\\
\fld{lmp} & $\lambda$ & — & — & 空 & 节点电价（等式行对偶恢复）\\
\fld{branch\_mu\_lower} / \fld{branch\_mu\_upper} & $\mu^-,\mu^+$ & — & — & 空 & 支路限值对偶（当前不可靠，见 \fld{branch\_mu\_valid}）\\
\fld{branch\_mu\_valid} & — & — & true/false & false & 恒为 false：原生支撑 LP 不能稳健恢复盒式变量对偶（字段注释 include/hacdcpf/optimal\_power\_flow/opf\_result.hpp:DCOPFResult；实现 src/optimal\_power\_flow/dc\_opf.cpp:extract\_dc\_opf\_result，设计说明 :621--634）\\
\fld{branch\_mu\_validity\_reason} & — & — & — & 空 & 上述判定的人读原因\\
\fld{lmp\_valid} & — & — & true/false & false & \fld{lmp} 是否来自报告原始解的等式行对偶（不蕴含阻塞对偶有效）\\
\fld{lmp\_validity\_reason} & — & — & — & 空 & 判定原因；\fld{compute\_lmp=false} 时注明被选项关闭（src/optimal\_power\_flow/dc\_opf.cpp:solve\_dc\_opf）\\
\fld{load\_shedding\_mw} & $\Delta P^d$ & MW & $\ge 0$ & 空 & 各母线切负荷量（src/optimal\_power\_flow/dc\_opf.cpp:extract\_dc\_opf\_result）\\
\fld{total\_load\_shedding\_mw} & — & MW & $\ge 0$ & 0 & 总切负荷量\\
\fld{solver\_chain} & — & — & — & 空 & 依次尝试的后端链，形如 ``\textless backend\textgreater:\textless status\textgreater''（字段注释 include/hacdcpf/optimal\_power\_flow/opf\_result.hpp:DCOPFResult；记录点 src/optimal\_power\_flow/dc\_opf.cpp:solve\_dc\_opf）\\
\fld{objective\_model} & — & — & QP/LP/LP-PWL & 空 & 上报目标实际使用的成本模型（实现 src/optimal\_power\_flow/dc\_opf.cpp:solve\_dc\_opf；头文件注释 :222--226 只列 ``QP''/``LP''，代码另发出 ``LP-PWL''，以代码为准）\\
\fld{pwl\_segments\_effective} & — & — & $\ge 0$ & 0 & LP 近似实际使用的线性段数；0 表示真 QP 成本\\
\fld{converter\_model\_scope} & — & — & — & 默认 & 恒为 ``dc-opf:linear-no-converter-model''，全部换流器标志为 false（字段注释 include/hacdcpf/optimal\_power\_flow/opf\_result.hpp:DCOPFResult；实现 src/optimal\_power\_flow/dc\_opf.cpp:extract\_dc\_opf\_result）\\
\fld{model\_limitations} & — & — & — & 空 & 能力边界声明（见下）\\
\end{paramtable}

\subsubsection*{诚实口径：model\_scope、ValidityFlags 与 model\_limitations}

平台约定每个 OPF 结果必须机器可读地声明自己\textbf{实际}建模了哪些换流器
物理、省略了哪些物理或证书，而不是让下游假定全保真。载体是
\fld{ConverterModelScope}
（\srcpath{include/hacdcpf/model/converter\_model\_scope.hpp:ConverterModelScope}）：
短标签 \fld{model\_scope} 加 12 个按特性划分的 \fld{ValidityFlags} 布尔位
（:ValidityFlags，覆盖 VSC 损耗、容量圆、电流限值、调制度、VDC 控制、多源协调、
DC/DC 损耗与占空比、LCC 准稳态与分接头、方程封闭性检查）。

AC OPF Parity 路径的标签\textbf{按本次运行实际强制的约束族拼接}——被
调用方关闭的约束族不出现在标签中
（\srcpath{src/optimal\_power\_flow/ac\_opf.cpp:solve\_with\_parity\_ipm}），例如
\fld{enforce\_converter\_capacity=false} 时标签不含
\fld{+capacity-circle}，对应 \fld{validity} 位同步置 false（:solve\_with\_parity\_ipm）。
传统 Native 路径标签固定为
\fld{ac-opf-native:vsc-free-pq+capacity-circle} 且
\fld{vsc\_vdc\_control\_modelled=false}（:solve\_ac\_opf\_impl）。DC OPF 标签固定为
\fld{dc-opf:linear-no-converter-model}（src/optimal\_power\_flow/dc\_opf.cpp:extract\_dc\_opf\_result），LCC 拒绝情形为
\fld{dc-opf:rejected-lcc}（:solve\_dc\_opf）。

\fld{model\_limitations} 是人读字符串列表，专门记录近似、省略的物理与不可
用证书（字段注释 include/hacdcpf/optimal\_power\_flow/opf\_result.hpp:ACOPFResult）。当前代码会写入的条目包括：

\begin{itemize}
  \item 非凸性声明：``AC OPF is a non-convex nonlinear program; convergence
        certifies a local KKT point, not a global optimum.''——所有 AC OPF
        收敛结果恒带（src/optimal\_power\_flow/ac\_opf.cpp:separate\_external\_grid\_dispatch（lambda））；
  \item 储能为单时段功率注入、不建模跨时段 SOC 动态（AC：:3207--3210；
        DC：src/optimal\_power\_flow/dc\_opf.cpp:extract\_dc\_opf\_result）；
  \item 能量路由器端口平衡为无损模型（src/optimal\_power\_flow/ac\_opf.cpp:separate\_external\_grid\_dispatch（lambda））；
  \item 传统 Native AC 路径经外部栅栏/罚环处理非线性工程限值，建议生产
        混合建模使用 ParityIPM（:separate\_external\_grid\_dispatch（lambda））；
  \item LCC 相关四条：功率/电流/$\alpha$/$\gamma$ 指令与换流变分接头为固定
        输入而非经济调度变量；额定电流饱和建模但 $\alpha$/$\gamma$ 限值为
        求解后审计而非硬约束；换相重叠迭代、分接头控制与主动损耗未建模；
        LCC 等式二阶导用解析雅可比的局部有限差分（:solve\_with\_parity\_ipm）；整流/
        逆变站求解后触发角/关断角窗口违例时逐站追加（:solve\_with\_parity\_ipm）；
  \item DC OPF 无损线性化、不建模电压幅值/无功/换流器物理
        （src/optimal\_power\_flow/dc\_opf.cpp:extract\_dc\_opf\_result）；LCC 拒绝情形（:solve\_dc\_opf）；
  \item LMP/阻塞对偶不可用时的原因说明（src/optimal\_power\_flow/ac\_opf.cpp:solve\_with\_parity\_ipm、
        3215--3218；src/optimal\_power\_flow/dc\_opf.cpp:extract\_dc\_opf\_result）。
\end{itemize}

\subsection{问题描述与求解器接口}\label{subsec:overview-interface}

\subsubsection*{OPFProblem：自包含任务描述}

\compmeta{OPFProblem}{\srcpath{include/hacdcpf/optimal\_power\_flow/opf\_problem.hpp:OPFProblem}}

\fld{OPFProblem} 把“优化什么”与“如何求解”解耦：一个非拥有系统指针
\fld{sys}、一份 \fld{ACOPFOptions}、一份 \fld{DCOPFOptions}，加
\fld{is\_ac}/\fld{is\_dc} 二选一标志。它由投影/调用层产生，可被任何
\fld{IOPFSolver} 实现消费。

\subsubsection*{IOPFSolver 与 DefaultOPFSolver}

\compmeta{IOPFSolver}{\srcpath{include/hacdcpf/optimal\_power\_flow/opf\_solver\_interface.hpp:IOPFSolver}}

求解器抽象只有一个纯虚方法
\texttt{solve(const OPFProblem\&) const}，返回类型
\fld{OPFSolveResult} 是 \texttt{std::variant<ACOPFResult, DCOPFResult>}
（:OPFSolveResult）。默认实现 \fld{DefaultOPFSolver}（:DefaultOPFSolver）不做第二套建模，直接
适配生产入口：\fld{sys} 为空或 \fld{is\_ac==is\_dc} 时抛
\fld{std::invalid\_argument}（:solve），否则按标志转发
\srcpath{hacdcpf::solve\_ac\_opf} 或 \srcpath{hacdcpf::solve\_dc\_opf}
（:solve）。该设计保证接口层与生产路径永远同源（头文件注释 :3--5）。

\subsubsection*{Parity 公共抽象层}

Parity 全空间建模的公共抽象（变量/约束布局与全部 KKT 求值函数）集中
在 \srcpath{include/hacdcpf/optimal\_power\_flow/formulation.hpp}，供
Parity IPM 与嵌入式 Ipopt 共用：

\begin{itemize}
  \item \fld{ParityOptions}（:ParityOptions）：建模族选项，含 \fld{load\_shedding}、
        \fld{voll}（0 表示按最大发电成本自动折算）、\fld{eps\_iac}
        （$|I_{ac}|$ 平滑）、目标与四个约束族开关；
  \item \fld{VarIndex}（:VarIndex）：全空间变量向量 $x$ 的连续分块布局
        （$\theta$、$V$、$P^g$、$Q^g$、$V^{dc}$、$P^{ac}_c$、$Q^{ac}_c$、
        $P^{dc}_c$、$\Delta P^d$、$\Delta Q^d$、$P^{ren}$、$Q^{ren}$、
        $P^{stor}$、$Q^{stor}$、$P^{stor,dc}$、$P^{dcdc}$、$P^{flex}$
        与能量路由器端口块），空块长度为零但保留有效偏移；
  \item \fld{ConstraintIndex}（:ConstraintIndex）：等式行（AC 有功/无功平衡、DC
        平衡、VSC/DC-DC/能量路由器耦合、AC 相角参考、DC 电压参考）与非线性
        不等式行（支路两端容量、VSC 容量圆、DC 支路、VSC 电流、调制度上下
        限、DC/DC 占空比）的连续行块布局；
  \item \fld{Problem}（:Problem）：装配完成的模型 + 全部模型到向量的映射
        （含 ZIP 系数、固定注入、$G/B$ 对角与 DC 电导矩阵）；
  \item 装配与求值函数：\srcpath{build\_problem}（:build\_problem）、
        \srcpath{build\_variable\_bounds}（:build\_variable\_bounds）、
        \srcpath{build\_initial\_point}（严格内点，:build\_initial\_point）、
        \srcpath{objective}（:objective）、
        \srcpath{objective\_gradient\_hessian\_diag}（:objective\_gradient\_hessian\_diag）、
        \srcpath{equality\_constraints}（:equality\_constraints）、
        \srcpath{equality\_jacobian}（:equality\_jacobian）、
        \srcpath{nonlinear\_inequality\_constraints}（:nonlinear\_inequality\_constraints）、
        \srcpath{nonlinear\_inequality\_jacobian}（:nonlinear\_inequality\_jacobian）、
        \srcpath{lagrangian\_hessian\_dense}/\srcpath{lagrangian\_hessian}
        （:lagrangian\_hessian）。数学口径详见第\ref{sec:parityform}节。
\end{itemize}

IPM 求解核的公共契约是 \fld{IPMOptions}/\fld{IPMResult} 与
\srcpath{solve\_primal\_dual\_ipm}
（\srcpath{include/hacdcpf/optimal\_power\_flow/native\_ipm\_solver.hpp:IPMOptions}、
\srcpath{include/hacdcpf/optimal\_power\_flow/native\_ipm\_solver.hpp:IPMResult}、
\srcpath{include/hacdcpf/optimal\_power\_flow/native\_ipm\_solver.hpp:solve\_primal\_dual\_ipm}），
另有 Davidenko 切向量接口 \srcpath{homotopy\_tangent}（:homotopy\_tangent）；详见
第\ref{sec:parityipm}节。高层驱动 \srcpath{solve\_ac\_opf\_parity}
（\srcpath{include/hacdcpf/optimal\_power\_flow/parity\_ipm\_solver.hpp:solve\_ac\_opf\_parity}）
串联装配、求解与结果回填，成功时置
\fld{solver\_path=ParityIPM}。雅可比对照诊断入口
\srcpath{compute\_jacobian\_diagnostics}
（\srcpath{include/hacdcpf/optimal\_power\_flow/ac\_opf\_solver.hpp:compute\_jacobian\_diagnostics}）
用于解析导数对有限差分的回归检查，见第\ref{sec:verification}节。

\subsection{嵌入式 Ipopt 后端：编译门控与运行时探测}\label{subsec:overview-ipopt}

嵌入式 Ipopt 是\textbf{可选}后端，其可用性由“编译期宏 + 运行时环境变量 +
能力查询”三层机制共同决定：

\begin{enumerate}
  \item \textbf{编译期门控}。CMake 选项 \srcpath{HACDCPF\_ENABLE\_IPOPT}
        仅 macOS 默认 ON、其余平台默认 OFF
        （\srcpath{CMakeLists.txt}:HACDCPF\_ENABLE\_IPOPT，帮助串为
        ``when supported by that checkout''）。非 Apple 平台显式开启
        \textbf{不再}触发 \fld{FATAL\_ERROR}：选项仅经
        \srcpath{MIPSOLVERS\_BUILD\_LOCAL\_IPOPT} 转发给 MIPSolvers
        子项目（同文件:MIPSOLVERS\_BUILD\_LOCAL\_IPOPT），能否真正构建
        取决于该 MIPSolvers 检出对内嵌 Ipopt 的支持；当前
        \srcpath{CMakeLists.txt} 中的 \fld{FATAL\_ERROR} 只剩 IPO 与
        测试联用、OpenDSS 库缺失等少数守卫。
        构建成功时 MIPSolvers 子项目向全部依赖目标传播宏
        \srcpath{HACDCPF\_HAVE\_IPOPT=1}
        （\srcpath{MIPSolvers/CMakeLists.txt:HACDCPF\_HAVE\_IPOPT}）。OPF 门面代码
        以 \texttt{\#ifdef HACDCPF\_HAVE\_IPOPT} 包裹引擎头文件包含
        （\srcpath{src/optimal\_power\_flow/ac\_opf.cpp:HACDCPF\_HAVE\_IPOPT}）与整个
        NLP 适配结构 \fld{ParityNLPInput}/\srcpath{build\_parity\_nlp\_input}
        （:ParityNLPInput、:build\_parity\_nlp\_input）。
  \item \textbf{运行时探测与降级}。统一入口
        \srcpath{solve\_parity\_with\_ipopt}
        （\srcpath{src/optimal\_power\_flow/ac\_opf.cpp:solve\_parity\_with\_ipopt}）以
        出参 \fld{available} 报告可用性：未编译时直接返回
        \fld{available=false}、\texttt{status=``Ipopt not compiled in''}
        （:solve\_parity\_with\_ipopt）；已编译时还可用环境变量
        \srcpath{HACDCPF\_DISABLE\_IPOPT\_OPF} 强制禁用（用于后端隔离与
        诊断，:solve\_parity\_with\_ipopt）。调用方按 \fld{available} 降级：显式
        \fld{Ipopt} 请求但不可用时回退原生 Parity IPM 并在后端标签注明
        ``(ipopt unavailable)''（:solve\_with\_parity\_ipm）；\fld{Auto} 模式先跑原生
        Parity IPM，不收敛且 \fld{allow\_fallback=true} 时尝试 Ipopt，
        成功则标记 ``ipopt\_filter\_linesearch (auto-fallback)''
        （:solve\_with\_parity\_ipm）。
  \item \textbf{能力查询}。\fld{SolverCapabilities::has\_ipopt} 在同一宏下
        置真（\srcpath{src/api/hacdcpf.cpp:get\_solver\_capabilities}），供 GUI/脚本在
        提交任务前探测（见第\ref{subsec:overview-capabilities}节）。
\end{enumerate}

Ipopt 路径还有两条必须知晓的诚实口径：其一，Ipopt 在大型混合模型上可能
以 \fld{Maximum\_Iterations\_Exceeded} 停机但已到达高度可用点，门面在原始
可行性严格满足、对偶/互补残差进入“可接受带”（对偶残差
$\le 1000\,\varepsilon_{\mathrm{dual}}$、互补残差
$\le \max(10^{-2},\,1000\,\varepsilon_{\mathrm{compl}})$，src/optimal\_power\_flow/ac\_opf.cpp:solve\_parity\_with\_ipopt）
时接受该点并在 \fld{status} 注明
``acceptable residuals after \ldots''（:solve\_parity\_with\_ipopt）；其二，
适配器不把约束乘子回传到 LMP 提取链路，Ipopt 解的
\fld{lmp\_valid=false}，原因为 ``The embedded Ipopt adapter does not
return constraint multipliers''（:solve\_with\_parity\_ipm）。Ipopt 终止时若对偶向量
维度齐全，门面会把约束/盒式对偶映射回 Parity 行序并做正性修复
（:solve\_parity\_with\_ipopt），供互补性度量使用，但这不改变 LMP 不可用的结论。

\subsection{OPF 在平台中的调用方速览}\label{subsec:overview-callers}

\begin{itemize}
  \item \textbf{时序生产模拟（UC $\to$ OPF $\to$ PF）}：
        \srcpath{src/time\_series/time\_series\_pf.cpp} 的三段流水线在
        UC 出力计划之后对每个时段调用 \srcpath{opf::solve\_ac\_opf}
        （:solve\_time\_series\_pf），上一时段收敛结果经 \fld{warm\_start} 传入下一
        时段（:solve\_time\_series\_pf），OPF 调度回写系统后再做牛顿 PF 校验；OPF 不
        收敛的时段回退为“UC 计划直接 PF”并如实记录（:solve\_time\_series\_pf）。
        按日并行分解时，各日分片强制使用线程安全的 ParityIPM 后端
        （:solve\_time\_series\_pf）。
  \item \textbf{电力市场（SCUC/SCED）}：
        \srcpath{src/market/market\_simulation.cpp} \textbf{不调用}本章
        的 \srcpath{solve\_dc\_opf}/\srcpath{solve\_ac\_opf}，而是在
        engine 层自建 SCUC/SCED 模型并用原生 B\&C/HiGHS/SCIP/Gurobi
        适配器求解（适配器选择 :803--839，SCED 装配与校验
        :2440--2469）。二者同属 DC 线性近似家族，但实现独立，价格/结算
        口径以市场模块自身结果为准。
  \item \textbf{网络重构}：
        \srcpath{src/network\_reconfiguration/topology\_reconfiguration.cpp}
        同样不经过 OPF 模块：其 MILP 自带虚拟潮流变量布局（文件头注释
        :9--17），经 engine 层 B\&C/MILP 适配器求解（:try\_native\_bc（lambda））。
  \item \textbf{其他内部调用方}：可靠性评估以 DC OPF 做削负荷最小化
        （\srcpath{src/reliability/reliability\_assessment.cpp:evaluate\_state}）；
        EV-交通耦合的联合社会福利与 CTM 模型内嵌 DC OPF
        （\srcpath{src/ev\_power\_traffic/joint\_social\_welfare.cpp:solve\_joint\_social\_welfare}、
        \srcpath{src/ev\_power\_traffic/ctm\_joint\_welfare.cpp:simulate\_ev\_power\_traffic\_ctm\_joint}）；
        反事实规划调用 AC OPF
        （\srcpath{src/analysis/counterfactual\_planning.cpp}\allowbreak :144）；
        SPPT 变形测试以 DC OPF 做投影前后对照
        （\srcpath{src/sppt/metamorphic.cpp}\allowbreak :223--225）。
  \item \textbf{HTTP 运行时 API（v1）}：GUI/脚本经
        \texttt{POST /api/v1/sessions/\{id\}/jobs} 提交
        \fld{analysis="optimal\_power\_flow"} 任务
        （路由 \srcpath{src/server/runtime\_api\_v1.cpp:register\_routes}；
        分派 :864--881）。请求体 \fld{solver} 字段映射为
        \fld{ac\_solver\_backend}（\fld{parity}/\fld{ipopt}/\fld{auto}/
        \fld{dispatch}，:ac\_opf\_options）或 DC OPF（\fld{solver="dc"}，
        :874--876）；\fld{network\_model} 当前仅支持
        \fld{balanced\_aggregate}（:execute\_analysis）。AC 结果序列化前自动执行
        \srcpath{verify\_opf\_result} 审计（:serialize\_ac\_opf），响应携带
        \fld{model\_scope}、\fld{validity\_flags}、
        \fld{model\_limitations}、\fld{infeasibility\_hints} 与
        \fld{audit}（:serialize\_ac\_opf）；DC 结果另携带 \fld{solver\_chain}、
        \fld{objective\_model} 与 \fld{branch\_mu\_valid}（:serialize\_dc\_opf）。
\end{itemize}

\subsection{求解器能力运行时查询}\label{subsec:overview-capabilities}

\compmeta{SolverCapabilities}{\srcpath{include/hacdcpf/api/solver\_capabilities.hpp:SolverCapabilities}}

\srcpath{hacdcpf::get\_solver\_capabilities()}（同头文件:44；实现
\srcpath{src/api/hacdcpf.cpp:get\_solver\_capabilities}）返回按编译结果填充的能力结构，
开销极小、可反复调用。恒真字段（内置 Eigen SparseLU、自研 IPM、AC/DC OPF、
三相潮流、二次目标）在:1983--1988 直接置真；可选后端按
\srcpath{HACDCPF\_HAVE\_*} 系列宏逐个置位（:get\_solver\_capabilities），这些宏由
MIPSolvers 子项目在找到对应库时传播（如
\srcpath{MIPSolvers/CMakeLists.txt:HACDCPF\_HAVE\_IPOPT} 对 Ipopt 的传播）。

\begin{paramtable}
\fld{has\_sparse\_lu} & — & — & true/false & true & Eigen SparseLU（恒有）\\
\fld{has\_klu} & — & — & true/false & false & SuiteSparse KLU（\texttt{\#ifdef HACDCPF\_HAVE\_KLU}，src/api/hacdcpf.cpp:get\_solver\_capabilities）\\
\fld{has\_umfpack} & — & — & true/false & false & SuiteSparse UMFPACK（:get\_solver\_capabilities）\\
\fld{has\_accelerate} & — & — & true/false & false & Apple Accelerate 框架（仅 macOS，:get\_solver\_capabilities）\\
\fld{has\_highs} & — & — & true/false & false & HiGHS LP/MIP；同时置 \fld{supports\_integer\_variables}（:get\_solver\_capabilities）\\
\fld{has\_gurobi} & — & — & true/false & false & Gurobi LP/QP/MIP（:get\_solver\_capabilities）\\
\fld{has\_ipopt} & — & — & true/false & false & 嵌入式 Ipopt NLP（:get\_solver\_capabilities；对应 CMake 门控见第\ref{subsec:overview-ipopt}节）\\
\fld{has\_scip} & — & — & true/false & false & SCIP MIP（:get\_solver\_capabilities）\\
\fld{has\_papilo} & — & — & true/false & false & PaPILO 预处理库（:get\_solver\_capabilities）\\
\fld{has\_native\_ipm} & — & — & true/false & true & 内置原始对偶 IPM（恒有）\\
\fld{has\_excel} & — & — & true/false & false & ETAP Excel I/O（\texttt{\#ifdef HACDCPF\_ENABLE\_ETAP}，:get\_solver\_capabilities）\\
\fld{has\_opendss} & — & — & true/false & false & OpenDSS 桥（\texttt{\#ifndef HACDCPF\_NO\_OPENDSS}，:get\_solver\_capabilities）\\
\fld{supports\_quadratic\_objective} & — & — & true/false & true & 二次目标（原生 + HiGHS 路径）\\
\fld{supports\_integer\_variables} & — & — & true/false & false & 整数变量（需 HiGHS/Gurobi/SCIP 任一）\\
\fld{supports\_ac\_opf} & — & — & true/false & true & 非线性 AC OPF（恒有）\\
\fld{supports\_dc\_opf} & — & — & true/false & true & 线性 DC OPF（恒有）\\
\fld{supports\_three\_phase} & — & — & true/false & true & 三相不平衡潮流（恒有）\\
\end{paramtable}

注意该查询只反映\textbf{编译进本构建}的后端，不代表运行时一定可用：
Gurobi 还需许可证（DC OPF 的 \srcpath{try\_gurobi\_qp} 每次调用都做
\fld{available()} 探测，src/optimal\_power\_flow/dc\_opf.cpp:try\_gurobi\_qp（lambda））；Ipopt 还可被环境变量
\srcpath{HACDCPF\_DISABLE\_IPOPT\_OPF} 临时禁用
（src/optimal\_power\_flow/ac\_opf.cpp:solve\_parity\_with\_ipopt）。因此生产代码应把能力查询作为提交前提示，
并以结果中的 \fld{solver\_name}/\fld{solver\_chain}/
\fld{profiling.linear\_solver\_backend} 作为实际所用后端的最终依据。

\subsection{本章小结与阅读路线}\label{subsec:overview-reading}

OPF 模块的公共契约可概括为：统一 Rich 输入、投影装配、按后端分派求解、
结果归因回稳定 ID，并以 \fld{model\_scope}/\fld{ValidityFlags}/
\fld{model\_limitations} 三元组机器可读地声明每次求解实际兑现的物理与
证书边界。后续章节按 AC OPF、DC OPF、交直流混合 OPF 三个主类与各专项
逐一展开：AC OPF（Native 紧致路径）见第\ref{sec:acopfnative}节，
DC OPF 见第\ref{sec:dcopf}节，交直流混合 OPF 的 Parity 建模层与 IPM 求解核
见第\ref{sec:parityform}、\ref{sec:parityipm}节，
RPO 见第\ref{sec:rpo}节，三相混合 OPF 见
第\ref{sec:threephase}节；AML 建模层见独立文档
\srcpath{docs/modules/power\_models/power\_models\_manual.tex}；跨引擎
验证与校核通道见第\ref{sec:verification}节。
